% titel: Blatt 8, A7: Hydrogensalze vs.\ Normalsalze löslich
% Teil II: der durchdachte Weg.

\subsection*{Gegeben und gesucht}
Erkläre, warum z.~B. NaHSO$_4$ in Wasser besser löslich ist als Na$_2$SO$_4$.

\subsection*{Ansatz}
Energiebilanz: Lösen gelingt gut, wenn
$\Delta H_{\text{Gitter}} - \Delta H_{\text{Hydratation}}$ klein ist.

\subsection*{Begründung}
\begin{itemize}
  \item \textbf{Gitterenergie.} In Na$_2$SO$_4$ steht ein zweifach geladenes
        Anion SO$_4^{2-}$ neben zwei Na$^+$. Der Coulomb-Term
        $q_1 q_2/r$ ist groß (doppelte Ladung!), das Gitter ist also sehr
        stabil und schwer aufzubrechen. NaHSO$_4$ enthält nur das einfach
        geladene HSO$_4^-$ $\Rightarrow$ deutlich kleinere Gitterenergie.
  \item \textbf{Hydratation.} Das HSO$_4^-$-Ion kann über sein Wasserstoff
        Wasserstoffbrücken zum Lösungsmittel aufbauen; auch ist es räumlich
        grösser und lässt sich gut von Wasser umgeben. Beides erhöht die
        Hydratationsenergie.
\end{itemize}
Damit ist für NaHSO$_4$ der Energieaufwand zum Lösen kleiner
(günstigere Bilanz) als für Na$_2$SO$_4$ $\Rightarrow$ bessere Löslichkeit.

\begin{ergebnis}
NaHSO$_4$ hat eine kleinere Gitterenergie (einfach geladenes Anion) und
bessere Hydratation durch H-Brücken als Na$_2$SO$_4$, daher höhere
Löslichkeit.
\end{ergebnis}

\begin{kontrolle}
Der Mechanismus erklärt auch NaHCO$_3$ vs.\ Na$_2$CO$_3$: das Hydrogen-Anion
löst sich im selben Muster besser. Ausnahmen existieren, der Trend ist aber
klar.
\end{kontrolle}

\begin{merksatz}
Die Löslichkeit eines Salzes entscheidet sich an der Bilanz aus
Gitterenergie (Coulomb, ~Ladungsprodukt) und Hydratationsenergie.
Einfach geladene, wasserstofftragende Anionen verschieben die Bilanz in
den gelösten Zustand.
\end{merksatz}
